\chapter{Methodology and Development}
\label{sec:methodik}

This chapter describes the conceptual development and technical implementation of the detection system presented in this thesis. It describes the path from the problem statement to the solution. In doing so, it examines in more detail the design decisions and the architecture of the individual components of the pipeline.

\section{Overall Architecture and Data Flow}
To reliably judge whether a cat is carrying a prey animal in its mouth, neither pure object detection nor pure image segmentation is sufficient. Since we only have a limited dataset of cats with prey animals available, it is essential to divide this task into several stages \cite{baumann2020smartcatflap}. This results in a pipeline consisting of object detection, object tracking, image segmentation and temporal state modelling. The final pipeline looks as follows:

\begin{enumerate}
	\item A frame is captured by the camera.
	\item The YOLO object detection analyses the frame.
	\begin{enumerate}
		\item If no cat was found, the pipeline ends here.
		\item If at least one cat is found, each cat face is extracted from the frame.
	\end{enumerate}
	\item The ByteTrack algorithm assigns an identity to the cat faces or creates a new one if the cat has only just entered the frame.
	\item The YOLO segmentation model examines each crop individually and returns for each cat a confidence score that it is carrying prey.
	\item The state machine aggregates the results of the image segmentation over time and then makes a final decision about which state (e.g.\ ``cat with prey'') the tracked objects are in.
	\begin{enumerate}
		\item If a cat is carrying prey in its mouth, the flap is closed.
		\item If all cats have been recognised as clean over the entire time window, the flap is opened.
	\end{enumerate}
\end{enumerate}

\section{Preprocessing for Infrared Night Vision}
Since cats do not only bring their prey home during the day under optimal lighting, a central requirement for the system is that it also works at night and in poor lighting. This brings a further difficulty: night vision cameras switch to infrared in the dark, which produces greyscale images.

So that the neural network does not learn the colour information, all images are converted to greyscale images before training and inference. Otherwise, neither object detection nor image segmentation would work reliably at night.

\section{Training and Adaptation of the Object Detector (YOLO26n)}
For object detection, the YOLO26 architecture in the \texttt{yolo26n} variant is used. It is the smallest one and offers high detection accuracy while providing the inference speed necessary for single-board computers.

Standard YOLO26 models are not specifically tuned to the detection of cat faces. Therefore, fine-tuning was carried out for exactly this problem \cite{yosinski2014transferable}. The basis for this is the Kaggle dataset (Cat Dataset by Crawford \cite{kaggle2024cat}), since it already contains annotations for the face. The dataset contains the coordinates for eyes, ears and mouth. To make these usable for YOLO, they have to be converted into a bounding box format. For this purpose, an algorithm was developed that determines the width of the box from the ears and the height from the ears and the mouth.

The fine-tuning of the YOLO26n model was carried out over 100 epochs at an image resolution of 640x640 pixels and with the standard YOLO data augmentation \cite{ultralytics_augmentation}. To avoid overfitting of the model, the training is stopped if no improvement occurs in the last ten epochs \cite{goodfellow2016deep}.

\section{Automated Generation of the Segmentation Dataset}
Unfortunately, there was no usable dataset for the segmentation of whether the cat is carrying prey in its mouth. Two kinds of images are required for the training: cats with and cats without prey. The cats without prey can again be taken from the Kaggle dataset (Cat Dataset by Crawford \cite{kaggle2024cat}). Cats with prey, on the other hand, are extremely difficult to obtain; they are either of poor quality or protected by copyright.

I therefore had the images of cats with prey generated by Gemini \cite{geminiteam2023gemini}. Using the official Gemini Application Programming Interface (API), I developed an algorithm that, taking the usage limit into account, automatically generates images of cats with prey in different poses and at different distances.

To prepare the dataset for fine-tuning the YOLO26n segmentation model, the previously trained cat face detector is used. It detects the cat faces in the images without prey and crops them in exactly the same way as later during inference. These crops are added to the dataset as negative examples without annotation. The images with prey were provided with segmentation masks of the prey animals using the AI-assisted annotation of Roboflow \cite{roboflow2025} and subsequently checked manually.

The segmentation model thus always has the same input conditions: it always receives the enlarged crop around the cat face without much background. This forms the perfect training basis, since during later inference the segmentation receives exactly the same input images.

\section{Crop Extraction}
Since the input of the segmentation model is the output of the detector, both must receive the required information. While the detector works best when it is trained only on the cat face, it also needs the information about what is located below the mouth of the cat face.

Therefore, the bounding box of the detector is modified before the segmentation model receives it. The cat face is extended downwards by twice the face height and sideways by 80\,\% of the face width on each side to make sure that the prey is in the image.

The detector expects a 3-channel image, so the greyscale image is duplicated three times (\texttt{cv2.cvtColor(gray, cv2.COLOR\_GRAY2BGR)}). The segmentation model, by contrast, was trained directly on single-channel greyscale images and receives the crop as a single-channel image. Beforehand, the crop is padded with black borders to form a square.

\section{Training of the Segmentation Model (YOLO Segmentation)}
As in the case of the detector, a YOLO model is used as the segmentation model. Since we are already in the Ultralytics environment and modern YOLO versions provide highly optimised networks for segmentation tasks, this decision was an obvious one.

The model receives an image of 640x640 pixels and returns a segmentation mask and a confidence score for each detected prey animal. Because the model's dataset is relatively small, strong emphasis is placed on data augmentation (random rotation of the image by up to 10°, horizontal flipping, variation of brightness and saturation). Mosaic augmentation, in which four images are combined into one, is deactivated so that small prey animals are not lost. These augmentation steps make the model more robust against unfamiliar poses of the cats.

\section{Object Tracking and Development of the State Model}
When a cat approaches the flap, the quality of the images changes considerably due to motion blur or unfavourable angles. As a result, the outputs of the segmentation model fluctuate. If, for example, the prey is briefly occluded, this would mean that the cat is allowed into the house.

To mitigate this problem, the object detection was extended with the ByteTrack algorithm. As a result, a cat receives a unique object ID when it first appears in the frame, and this ID remains the same for that cat in the following frames. Based on the uniqueness of each cat, a state machine could now be developed. This class \texttt{StateMachine} manages a \texttt{TrackState} object for each detected cat (\texttt{track\_id}). The model knows four states:

\begin{itemize}
    \item \texttt{NO\_CAT}: There is no relevant object in the field of view.
    \item \texttt{GATHERING\_HISTORY}: A cat has been detected, but there are not yet enough frames for a decision.
    \item \texttt{CAT\_NO\_PREY}: A cat has been detected but is not carrying prey.
    \item \texttt{CAT\_WITH\_PREY}: A cat has been detected and is carrying prey with sufficient certainty.
\end{itemize}

So that the system's decision is not made on the basis of the most recent frame alone, the \texttt{TrackState} object stores a history of the confidence scores of the segmentation model for the last 20 frames. The current state is recalculated for every frame, simply by computing the moving average of the history. If the threshold of 0.3 is reached or exceeded, the cat is marked as \texttt{CAT\_WITH\_PREY}. If the value falls below it, the cat is marked as \texttt{CAT\_NO\_PREY}, provided that 20 values have already been collected.

The threshold of 0.3 was chosen on the basis of the parameter series in Section~\ref{sec:ablation_sm}. Since the average over 20 frames already smooths out individual outliers considerably, the threshold can be set low without the cat having to stay outside unjustly. However, this value is not set in stone. If it were more important to the end user that the cat is allowed in when in doubt, the value would be set higher, for example at 0.5.

The decision about the flap results from the states of all cats. If at least one cat is carrying prey, the flap is locked immediately. It is only unlocked when all cats in the frame have a complete history and are considered clean. Only cats that are visible in the current frame are taken into account. If the tracker assigns a new identity to a cat, the short-lived track left behind therefore does not block the release. A cat with prey, on the other hand, maintains the lock until it is removed from memory. So that the flap does not constantly open and close, a renewed release is permitted at the earliest five seconds after the last change.

To prevent memory leaks in the state machine, cats that have left the field of view or that have not been found for more than 30 frames (\texttt{max\_missed\_frames}) are removed from memory.
